\section{Introduction}
\label{sec:introduction}

Autonomous language agents are increasingly deployed to execute long-horizon, multi-step problem solving across software engineering, scientific discovery, and automated data workflows~\citep{pushkin2026lead,zhao2026scienceflow,afoud2026vreact}. As task horizons expand from single-turn tool calls to extended graphs consisting of dozens of interdependent operations, autonomous agents exhibit severe planning-quality degradation~\citep{wang2026reasoning,chen2026horizon}. In unverified sequential execution, errors introduced early in an execution trajectory propagate downstream, corrupting intermediate datasets, invalidating subsequent inferences, and ultimately resulting in catastrophic task failures~\citep{sinha2025illusion,venkataramani2026masprove}.

To combat compounding error propagation, recent literature emphasizes runtime verification and execution rollback~\citep{zhuang2026agentrewind,yu2026rir}. Rather than relying solely on post-hoc evaluation after complete workflow execution (terminal verification), autonomous systems can evaluate intermediate execution states against formal semantic or structural specifications, rolling back corrupted state variables to previously validated checkpoints when anomalies are detected~\citep{shoham2026jevact,afoud2026vreact}.

However, intermediate verification introduces a critical engineering and resource trade-off. Formal verification calls—whether implemented via deterministic relational validators, schema linters, or auxiliary language-model judges—incur non-trivial computational latency, memory overhead, and monetary invocation costs~\citep{venkataramani2026masprove,yang2025agenticrobot}. Verifying the agent state at every execution step is computationally prohibitive in large-scale workflows, whereas deferring verification entirely until workflow termination maximizes error propagation, leading to substantial wasted computation and large rollback distances~\citep{pushkin2026lead,zhuang2026agentrewind}.

Consequently, autonomous workflow systems must operate under a \emph{verification budget constraint}: given a finite allocation of verification opportunities $K$ across a workflow of horizon $H$ (where $K \ll H$), \textbf{where should verification checkpoints be placed?}

Two prominent scheduling paradigms dominate current architectural discussions:
\begin{enumerate}
    \item \textbf{Fixed-Periodic Scheduling:} Verification is invoked at uniform temporal intervals (e.g., every $\lceil H/3 \rceil$ steps). This clock-like schedule is oblivious to internal task graph topology but guarantees a strict upper bound on the maximum number of steps an unmonitored error can travel~\citep{yang2025agenticrobot}.
    \item \textbf{Dependency-Based Milestone Scheduling:} Verification checkpoints are placed at critical topological boundaries, specifically immediately prior to multi-parent dependency integrations (join nodes). Intuitive architectural principles suggest that validating intermediate artifacts before they merge into downstream dataflows should prevent cross-branch contamination~\citep{afoud2026vreact}.
\end{enumerate}

Despite intuitive arguments favoring topological milestone placement, existing literature lacks a rigorously controlled empirical comparison between periodic and milestone-based verification under \textbf{strictly matched verification-call budgets}. Existing studies either evaluate periodic schedules in isolation~\citep{yang2025agenticrobot}, propose integration-level verifiers without experimental schedule pilots~\citep{afoud2026vreact}, or evaluate unconstrained process-reward judges that conflate check frequency with check placement~\citep{venkataramani2026masprove}.

\paragraph{Contributions of this Work}
In this study, we conduct an extensive, pre-registered empirical investigation into verification scheduling policies within long-horizon directed acyclic graph (DAG) workflows. Across 90 synthetic task graph clusters evaluating 4,680 execution episodes under controlled fault injection, we systematically evaluate four verification arms:
\begin{itemize}
    \item \textbf{Arm A (Final-Only):} Verification deferred strictly until terminal step $H$ ($K_{\text{alloc}} = 1$).
    \item \textbf{Arm B (Fixed-Periodic Baseline):} Verification scheduled at regular intervals $\lceil H/3 \rceil, \lceil 2H/3 \rceil, H$ ($K_{\text{alloc}} = 3$).
    \item \textbf{Arm C (Dependency-Milestone Intervention):} Verification scheduled immediately prior to Join 1, Join 2, and terminal step $H$ ($K_{\text{alloc}} = 3$).
    \item \textbf{Arm D (Every-Step Bound):} Verification invoked after every execution step ($K_{\text{alloc}} = H$).
\end{itemize}

By enforcing an exact matched budget between Arms B and C ($K_{\text{alloc}} = 3$, realized mean $2.7436 \pm 0.5872$ calls across both arms, $p = 1.0000$), we isolate the pure effect of checkpoint placement. Our central empirical findings reveal:
\begin{enumerate}
    \item \textbf{Empirical Rejection of Milestone Superiority:} Contrary to initial theoretical intuition, milestone-based verification underperforms fixed-periodic verification on downstream error containment. Milestone verification incurs a statistically significant $3.91\%$ increase in the downstream taint-consumption fraction ($\Delta \text{TCF} = +0.0391$, paired permutation $p < 0.0001$, 95\% bootstrap CI $[+0.0366, +0.0414]$, Cohen's $d = 3.35$).
    \item \textbf{Detection Latency Divergence:} Fixed-periodic verification catches contract-violating faults with a mean delay of $5.50$ execution steps, whereas milestone verification requires $7.81$ steps ($\Delta \text{Delay} = +2.31$ steps, $p < 0.0001$).
    \item \textbf{Topological Clumping Mechanism:} We demonstrate that dependency integration nodes naturally cluster in the middle and late stages of multi-branch DAGs. This clumping leaves early computational branches unmonitored under milestone schedules, allowing early corruptions to cascade across multiple nodes before reaching a join boundary.
    \item \textbf{Containment vs. Completion Decoupling:} Under a single-replay recovery policy, intermediate verification does not alter eventual workflow success rates ($38.46\%$ across all arms), but significantly reduces recovery re-execution waste, saving up to $3.81$ re-executed nodes per episode.
\end{enumerate}
